\section{Discussion}
\label{sec:discussion}

\subsection{Limitations}

\todo{Synthesize limitations surfaced across Sections~\ref{sec:model}
(partial-status programs and their specific gaps),
\ref{sec:data} (the \efrs{} construction is itself a modelled input, not an
observed one, for the synthetic high-income rows), and
\ref{sec:validation} (the two named-but-unassembled validation types:
ready reckoners and HBAI distributional statistics; the weaker
individual-level fit for student loan repayments; and the fact that the
\ukmod{} comparison in \ref{sec:validation-ukmod} compares against
\policyengine{}'s own calibration target for the ``Official'' series, not
an independent one, for the aggregate figures -- only the \ukmod{} column
is a genuinely external comparator there). State each limitation plainly
rather than burying it, consistent with the paper's stated policy-neutral,
facts-only tone.}

\subsection{Relationship to \ukmod{} and \euromod{}}

\policyengine{} UK and \ukmod{} are independently developed, open models of
overlapping scope; both descend from a broader European tradition of
survey-based tax-benefit microsimulation established by \euromod{}
\citep{sutherland2013euromod}. This paper's validation-comparative sections
report where the two models' calculated aggregates agree or differ from
published or cited figures; \todo{add a sentence noting any structural
differences in scope or approach that are already publicly documented (e.g.
policyengine-uk\#300's list of differences, if that issue's content has been
resolved into a citable public document) -- do not introduce new
comparative claims not already public.} Consistent with the collegial
framing required for this paper, this section makes no claim about which
model is preferable for a given use case beyond what a reader can verify
from the cited figures themselves.

\subsection{Reproducibility}

\todo{State how a reader can reproduce every table in this paper: which
\texttt{policyengine-uk} and \texttt{populace} commits/tags were used, and
whether (following the sibling imputation paper's disclosure convention)
this paper's own repository contains scripts that regenerate each table
from those pinned versions, or whether tables are transcribed by hand from
notebook output with a stated commit reference. The imputation-paper
precedent regenerates every table from committed run configurations; decide
whether this paper's validation notebooks (currently living in
policyengine-uk, not in this paper's own repo) should be re-run by CI here,
mirrored into this repo, or cited by pinned commit only, and record that
decision.}

\subsection{Future versions}

\todo{Note known near-term changes that would trigger a v2 per the
versioning framing in Section~\ref{sec:introduction} (e.g. a
\texttt{policyengine\_core} engine migration, or the Axiom rules-encoding
migration referenced in the wider PolicyEngine program, if and when it
reaches policyengine-uk) -- state only changes that are already planned and
documented elsewhere, not speculative ones.}
